\documentclass[10pt,a4paper]{amsart}
\usepackage[margin=19mm]{geometry}
\usepackage{amsmath,amssymb,mathtools}
\usepackage[scr=rsfs,cal=euler]{mathalfa}
\usepackage{xfrac}
\usepackage[unicode,hidelinks,bookmarks=false]{hyperref}
\newcommand{\ie}{i.e.\ }
\newcommand{\cf}{cf.\ }
\newcommand{\resp}{respectively\ }
\newcommand{\Subsection}{\subsection}
\providecommand{\nobreakdash}{\nobreak\hbox{-}}
\setlength{\emergencystretch}{2em}
\multlinegap0pt
\usepackage{amssymb}
\usepackage{mathtools}
\usepackage{xcolor}
\newtheorem{theorem}{Theorem}
\newtheorem{lemma}{Lemma}
\newcommand{\R}{\mathbb R}
\newcommand{\cL}{\mathcal L}
\newcommand{\supp}{\operatorname{supp}}
\title{Distribution of simplices in the discrete and continuous settings}
\author{Thang Pham, Chun-Yen Shen, Boqing Xue}
\date{Source: arXiv:2608.01274v1 (CC BY 4.0)}
\begin{document}
\maketitle

\noindent\emph{Source and licence.} Distribution of simplices in the discrete and continuous settings. Version arXiv:2608.01274v1 (CC BY 4.0), distributed by arXiv under the Creative Commons Attribution 4.0 International licence, http://creativecommons.org/licenses/by/4.0/, as recorded in the arXiv licence metadata for this exact version and shown on the abstract page https://arxiv.org/abs/2608.01274v1. Full licence notice: \texttt{licenses/arxiv-2608.01274-CC-BY-4.0.txt}.\par\smallskip

\noindent\textit{Editorial scope.} Counted unit: the article's theorem thm:euclidean-codim-one (source0.tex lines 273--290). The article's complete original proof of it is delivered with it (source0.tex lines 2834--3085, one \texttt{proof} environment of 252 lines, which the article places away from the statement). No mathematics is written, repaired or re-derived: the statement and the proof are the article's own lines, in the article's own notation, and no retained line is substituted or re-flowed.  Retained uncounted as the source-local context the proof needs: lines 1351--1354; lines 1414--1429; lines 2607--2617.  Nothing was omitted from any retained span, so the package carries no omission record.  No retained line contains a citation, so the wrapper carries no bibliography and no bibliography entry.  No retained line contains a figure environment, an image-inclusion command or drawing code, so no image is delivered and the package declares no asset dependency.  Editorial formatting only, in addition: an AMS \texttt{amsart} preamble that restates the article's own declarations and macros used by the retained lines, namely the environments \texttt{theorem}, \texttt{lemma} on separate counters, together with the standard packages those lines need.  The retained environments print in this document as Theorem 1, Lemma 1 in order of appearance, whereas the article prints the same items with the numbering of its own full document.  The complete native source of the article as distributed in the arXiv e-print for this exact version is bundled unchanged as \texttt{sources/206611/source0.tex}.
\begin{equation}
        \mu(B(x,r))\le C_\mu r^s.
\label{eq:frostman}
\end{equation}

\begin{lemma}
\label{lem:squaring-ac}
Let \(\Omega\) be a finite Borel measure supported in a bounded subset
of \([0,\infty)^M\), and let
\[
        \Theta=\operatorname{Sq}_\#\Omega.
\]
If \(\Theta\ll\cL^M\), then \(\Omega\ll\cL^M\).

Moreover, for \(1\leq j<d\), the distance tuples and the
squared-distance tuples of degenerate \(j\)-simplices form
Lebesgue-null subsets of \(\R^{t_j}\), where
\[
        t_j=\binom{j+1}{2}.
\]
\end{lemma}

\begin{lemma}[]
\label{lem:codim-marginal-descent}
For every $1\le j\le n$, there exists a Borel set
\[
G_j\subset (\operatorname{supp}\mu)^j,
\qquad
\mu^j(G_j)=1,
\]
such that $\rho_Y^{(j)}$ has an $L^2(\mathbb R^j)$ density for every
$Y\in G_j$.
\end{lemma}

\begin{theorem}\label{thm:euclidean-codim-one}
Let \(d\geq3\), let \(m=\binom d2\), and let \(E\subset\R^d\) be
compact.  If
\[
        \dim_{\mathrm H}(E)>d-1,
\]
then there exist a probability measure \(\mu\), supported on \(E\),
and a Borel set \(E_\mu\subset\supp\mu\), with \(\mu(E_\mu)=1\), such
that, for every \(x\in E_\mu\),
\[
        (\Phi_{d-1,x})_\#\mu^{d-1}\ll\cL^m.
\]
Consequently,
\[
        \cL^m\bigl(\Delta_{d-1,x}^{\mathrm{nd}}(E)\bigr)>0
        \qquad (x\in E_\mu).
\]
\end{theorem}

\begin{proof}[Proof of Theorem~\ref{thm:euclidean-codim-one}]
Let \(\mu\) be the \(s\)-Frostman probability measure fixed in
\eqref{eq:frostman}, where
\[
d-1<s<\dim_{\mathrm H}(E).
\]

For \(j=1\), interpret \(\mu^0\) as unit mass on a one-point space.
For every \(1\leq j\leq n\), define
\[
E_j
:=
\left\{
x\in\supp\mu:
\mu^{j-1}
\left(
\left\{
(z_1,\ldots,z_{j-1}):
(x,z_1,\ldots,z_{j-1})\notin G_j
\right\}
\right)
=0
\right\}.
\]
Since \(G_j\) is Borel, the corresponding section-measure function is
Borel, and hence \(E_j\) is Borel. By Fubini's theorem and
Lemma~\ref{lem:codim-marginal-descent},
\[
\mu(E_j)=1.
\]
Thus, for every \(x\in E_j\),
\begin{equation}
\mu^{j-1}
\left(
\left\{
(z_1,\ldots,z_{j-1}):
(x,z_1,\ldots,z_{j-1})\notin G_j
\right\}
\right)
=0.
\label{eq:codim-good-fixed-pin}
\end{equation}

Set
\[
E_\mu
:=
\bigcap_{j=1}^n E_j.
\]
Then \(E_\mu\subset\supp\mu\) is Borel and
\[
\mu(E_\mu)=1.
\]

Fix \(x\in E_\mu\), and write
\[
\lambda_{j,x}
:=
(\Phi_{j,x})_\#\mu^j.
\]
We prove by induction on \(j\) that
\[
\lambda_{j,x}\ll\cL^{t_j},
\qquad 1\leq j\leq n.
\]

For \(j=1\), condition~\eqref{eq:codim-good-fixed-pin} says that
\[
(x)\in G_1.
\]
Hence, by the defining property of \(G_1\),
\[
\lambda_{1,x}
=
\left(
z\longmapsto |z-x|^2
\right)_\#\mu
=
\rho^{(1)}_{(x)}
\ll\cL^1.
\]

Now suppose that \(2\leq j\leq n\) and
\[
\lambda_{j-1,x}\ll\cL^{t_{j-1}}.
\]
For
\[
Z=(z_1,\ldots,z_{j-1}),
\]
define
\[
\rho_{x,Z}
:=
\left(
z\longmapsto
\bigl(
|z-x|^2,
|z-z_1|^2,
\ldots,
|z-z_{j-1}|^2
\bigr)
\right)_\#\mu.
\]
By \eqref{eq:codim-good-fixed-pin},
\[
(x,Z)\in G_j
\]
for \(\mu^{j-1}\)-almost every \(Z\). Therefore, by the defining
property of \(G_j\),
\[
\rho_{x,Z}\ll\cL^j
\]
for \(\mu^{j-1}\)-almost every \(Z\).

After a fixed permutation of coordinates, identify
\[
\R^{t_j}
=
\R^{t_{j-1}}\times\R^j,
\qquad
t_j=t_{j-1}+j.
\]
Let
\[
N\subset\R^{t_{j-1}}\times\R^j
\]
be a Borel set of Lebesgue measure zero. For
\(b\in\R^{t_{j-1}}\), write
\[
N_b
:=
\{a\in\R^j:(b,a)\in N\}.
\]
By Fubini's theorem,
\[
\cL^j(N_b)=0
\]
for \(\cL^{t_{j-1}}\)-almost every \(b\). Since
\[
(\Phi_{j-1,x})_\#\mu^{j-1}
=
\lambda_{j-1,x}
\ll\cL^{t_{j-1}},
\]
it follows that
\[
\cL^j
\left(
N_{\Phi_{j-1,x}(Z)}
\right)
=0
\]
for \(\mu^{j-1}\)-almost every \(Z\).

For \(\mu^{j-1}\)-almost every such \(Z\), we also have
\(\rho_{x,Z}\ll\cL^j\). Hence
\[
\rho_{x,Z}
\left(
N_{\Phi_{j-1,x}(Z)}
\right)
=0.
\]
Integrating first in the new vertex gives
\[
\begin{aligned}
\lambda_{j,x}(N)
&=
\int
\rho_{x,Z}
\left(
N_{\Phi_{j-1,x}(Z)}
\right)
\,d\mu^{j-1}(Z)\\
&=0.
\end{aligned}
\]
Therefore
\[
\lambda_{j,x}\ll\cL^{t_j}.
\]
This completes the induction.

Taking \(j=n=d-1\) and recalling that \(t_n=m\), we obtain
\[
(\Phi_{d-1,x})_\#\mu^{d-1}
\ll\cL^m
\qquad
(x\in E_\mu).
\]

Let \(\Omega_x\) be the corresponding pinned ordinary-distance
configuration measure. Since \(E\) is compact, \(\Omega_x\) is
supported in a bounded subset of \([0,\infty)^m\), and
\[
(\Phi_{d-1,x})_\#\mu^{d-1}
=
\operatorname{Sq}_\#\Omega_x.
\]
Lemma~\ref{lem:squaring-ac} therefore gives
\[
\Omega_x\ll\cL^m.
\]

The same lemma shows that the ordinary-distance tuples of degenerate
\((d-1)\)-simplices form a Lebesgue-null set. Since
\(\Delta_{d-1,x}(E)\) is compact and the degenerate configurations are
characterized by the vanishing of the Gram-determinant polynomial
\[
\widetilde P_{d-1}(r)
:=
P_{d-1}
\bigl(
(r_{ab}^2)_{0\leq a<b\leq d-1}
\bigr),
\]
the set
\[
\Delta^{\mathrm{nd}}_{d-1,x}(E)
=
\Delta_{d-1,x}(E)
\cap
\left\{
r\in\R^m:
\widetilde P_{d-1}(r)\neq0
\right\}
\]
is Borel.

Moreover, the degenerate locus has \(\cL^m\)-measure zero, so
\(\Omega_x\) assigns zero mass to that locus. Since \(\Omega_x\) is a
probability measure,
\[
\Omega_x
\bigl(
\Delta^{\mathrm{nd}}_{d-1,x}(E)
\bigr)
=
1.
\]
Finally, since \(\Omega_x\ll\cL^m\),
\[
\cL^m
\bigl(
\Delta^{\mathrm{nd}}_{d-1,x}(E)
\bigr)
>0
\qquad
(x\in E_\mu).
\]
\end{proof}

\end{document}
